%% Condition 3 Bounded Process AI Declaration -- WIL Project 84 (group submission)
%% Built on the course template (_SCT_Generative_AI_Framework__Condition_3_Bounded_Use_Declaration_Form).
\documentclass[manuscript,screen,nonacm]{acmart}

\settopmatter{printacmref=false} % Removes citation information below abstract
\renewcommand\footnotetextcopyrightpermission[1]{} % removes footnote with conference information in first column
\pagestyle{plain} % removes running headers

\AtBeginDocument{%
  \providecommand\BibTeX{{%
    \normalfont B\kern-0.5em{\scshape i\kern-0.25em b}\kern-0.8em\TeX}}}

%Add any additional packages and commands here

%%
%% This command processes the author and affiliation and title
%% information and builds the first part of the formatted document.

\usepackage{hyperref}
\usepackage[most]{tcolorbox}

\usepackage{array}
\usepackage{enumitem}
\usepackage{geometry}
\newcommand{\wordcount}[1]{\noindent\textbf{[Word count: #1]}} %macro for word count

\newcommand{\todo}[1]{\textcolor{red}{\bf [[TODO: #1]]}} 
\newcommand{\instructions}[1]{\textcolor{blue}{\textbf{[[Remove before submitting:} #1]]}} %Instructions to be removed before submitting

\newcommand{\update}[1]{\textcolor{blue}{\textbf{[[Update before submitting:} #1]]}} %Instructions to be removed before submitting


%\newcommand{\instructions}[1]{} %Instructions to be removed before submitting

%extra packages for the tables in the appendix
\usepackage{longtable}
\usepackage{array}
\renewcommand{\arraystretch}{1.5}
\begin{document}

\title{AI Use Declaration -- Condition 3: Bounded Process AI}
\thanks{This declaration is part of your assessed submission.}

\begin{abstract}
\begin{tcolorbox}[colback=blue!05,colframe=black,]
\emph{Required for all Condition 3 assessments -- not graded, but reviewed.}
This document allows your marker to judge your authorship, your decision-making, and the quality of your critical engagement with AI. It is not used to reward AI use. It is used to assess whether you remained the author and final decision-maker. A concise, honest, and specific account is more useful than a long but vague one.
\end{tcolorbox}
\end{abstract}

\author{Arulmozhi Ezhilarasi Arunachalam (s4138509)}
\affiliation{\institution{RMIT University}\city{Melbourne}\country{Australia}}
\email{s4138509@student.rmit.edu.au}
\author{Dharani Suresh (s4118043)}
\affiliation{\institution{RMIT University}\city{Melbourne}\country{Australia}}
\email{s4118043@student.rmit.edu.au}
\author{Lingkesh Natanasabapathy (s4184049)}
\affiliation{\institution{RMIT University}\city{Melbourne}\country{Australia}}
\email{s4184049@student.rmit.edu.au}
\author{Prakash Sudarsanraj (s4135999)}
\affiliation{\institution{RMIT University}\city{Melbourne}\country{Australia}}
\email{s4135999@student.rmit.edu.au}
\author{Sakthisree Gaddam (s4136157)}
\affiliation{\institution{RMIT University}\city{Melbourne}\country{Australia}}
\email{s4136157@student.rmit.edu.au}
\renewcommand{\shortauthors}{WIL Project 84}

\maketitle

\section*{Part A: Assessment Details}

\noindent\textbf{Assessment number \& title:} WIL Project: Written Report (COSC2669/COSC2816 Case Studies in Data Science), group WIL Project 84, project \emph{SkillGap}.

\vspace{1em}
\noindent\textbf{Date submitted:} \hspace{1em} \rule{4cm}{0.4pt}

\section*{Part B: Account of AI Use}

\noindent Complete all five fields below for each AI Tool you used. If you used more than one tool for different purposes, copy this section for each. If you did not use any AI tool, select the statement below and you do not need to complete the remaining fields in Parts B nor C (go to Part D).

\begin{itemize}
 \item \noindent\CheckBox[name=declNotUse,bordercolor={0 0 0},width=1.2ex,height=1.2ex]{}~ I did not use any AI tool in the preparation or submission of this assessment.
 \item \noindent\(\boxtimes\)~I have used AI for the purposes approved for this assessment. Here are the details of my use.
\end{itemize}

\subsection*{AI\_TOOL\_1}
\begin{enumerate}
    \item \textbf{Name of the AI tool(s) used:}

    \vspace{0.5em}
    \noindent Claude Code (Anthropic), an AI coding assistant, using Claude Opus models.

    \vspace{1.5em}
    \item \textbf{What specifically did you use the AI tool for?}

    \vspace{0.5em}
    \begin{itemize}
      \item \emph{Code:} reviewing and fixing defects in the starter prototype, and writing parts of the retrieval, generation, evaluation and data-cleaning code (all changes are logged in \texttt{CHANGELOG.md}).
      \item \emph{Quality assurance:} an independent test audit (test plan, 72 test cases, bug report) on 5~October 2026.
      \item \emph{Annotation:} judging faithfulness and correctness of 54 generated answers, auditing 72 relevance labels, and hand-labelling skills in 20 ads. Each judgement and its reason are recorded in the annotation sheets.
      \item \emph{Report:} drafting the written report from the team's results, Milestone~1 document and change log; porting it into the course Overleaf template; checking every number against the result files; drafting the individual self-reflections from the project record.
      \item \emph{Contribution sheet:} filling in names, IDs, the agreed 20\% shares and work descriptions based on the Milestone~1 roles.
    \end{itemize}

    \vspace{1.5em}
    \item \textbf{Confirm the scope: what did the AI tool NOT do?}

    \vspace{0.5em}
    \noindent It did not choose the project domain, the stakeholders, the user story, the team roles or the equal contribution split; these were the team's decisions (Milestone~1 and team meetings). The Walert reproduction reported at Milestone~1 was done by the team. The AI did not create or alter any data: every number in the report comes from the repository's scripts and result files, and the AI-assisted annotation is labelled as such in the report, which recommends a human re-check. It did not access the team's Overleaf, GitHub or Trello accounts and did not submit anything. The team reviewed the outputs and made the final decisions on what to submit.
\end{enumerate}

\section*{Part C: Before-and-After Sample} (if applicable)
\begin{tcolorbox}[colback=blue!05,colframe=black,]
If the AI changed any text or code directly -- for example, if you accepted a revised sentence, corrected code, or edited paragraph -- attach a brief before-and-after sample below. This is required where AI directly altered submitted text or code. It is not required if AI was used only for brainstorming or outlining that you then developed independently.
\end{tcolorbox}

\subsection*{SAMPLE 1: project scope}

\noindent\textbf{Before (original text, from the team's Milestone~1 document):}

\vspace{0.5em}
\noindent\emph{``To develop and evaluate a Test-Driven RAG assistant that helps international data and IT students understand Australian graduate job requirements and identify evidence-based skill gaps for their target roles.''} \ldots\ \emph{``These risks will be managed using a fixed dataset snapshot, explicit evaluation criteria, a limited set of target roles, and preference for no-cost or locally executable tools where practical.''}

\vspace{1.5em}
\noindent\textbf{After (text as submitted, report Section~1.4):}

\vspace{0.5em}
\noindent\emph{``From `data and IT' to three data roles. Data Analyst, Business/Data Analyst and Data Scientist ads share a clear, testable skill vocabulary (SQL, Excel, Python, \ldots), so labels could be derived and audited reliably. \ldots\ This follows the Milestone~1 risk plan to limit the set of target roles.''} and \emph{``Only 21 of the 159 ads (13\%) mention work rights, citizenship, visas or sponsorship at all, and none mentions international students. SkillGap therefore keeps international students as a key user group. It answers work-rights questions from those ads \ldots\ and declines them otherwise, rather than guessing.''}

\vspace{1.5em}
\noindent\textbf{Annotation: what changed, and what did you decide independently?}

\vspace{0.5em}
\noindent The AI wrote the scope-change paragraph and measured the 21-of-159 figure from the knowledge base to justify it. The aim, the international-student user group, the fixed dataset snapshot, the local no-cost tools and the decision to limit target roles came from the team's Milestone~1 plan. The paragraph explains how those team decisions were carried through.

\section*{Part D: Declaration}

\noindent Select all that apply and sign below.

\begin{itemize}
\item \noindent\CheckBox[name=declNotUsed2,bordercolor={0 0 0},width=1.2ex,height=1.2ex]{}~I did not use any AI tool in the preparation or submission of this assessment. \\[0.8em]
\item \noindent\CheckBox[name=declUsePermitted,bordercolor={0 0 0},width=1.2ex,height=1.2ex]{}~I used AI only for the process purposes explicitly permitted in the task instructions.\\[0.8em]
\item \noindent\CheckBox[name=declRemainedAuthor,bordercolor={0 0 0},width=1.2ex,height=1.2ex]{}~I remained the author and final decision-maker throughout this assessment.\\[0.8em]
\item \noindent\CheckBox[name=declBeforeAfter,bordercolor={0 0 0},width=1.2ex,height=1.2ex]{}~The account I have provided above accurately and honestly describes how I used AI in this task.\\[0.8em]
\item \noindent\CheckBox[name=declResponsibility,bordercolor={0 0 0},width=1.2ex,height=1.2ex]{}~I take responsibility for the quality, accuracy, integrity, and final form of the submission.\\[0.8em]
\item \noindent\CheckBox[name=declInterview,bordercolor={0 0 0},width=1.2ex,height=1.2ex]{}~I understand that my assessor may ask me to explain or discuss my work and my AI use, and that I am expected to be able to do so without reference to AI-generated material.\\[0.8em]
\item \noindent\CheckBox[name=declIntegrity,bordercolor={0 0 0},width=1.2ex,height=1.2ex]{}~I understand that using AI beyond the permitted scope, or failing to declare AI use where required, may result in a mark reduction or referral for academic integrity review.
\end{itemize}

\vspace{2em}
\noindent\textbf{Signature:} \hrulefill

\vspace{2em}
\noindent\textbf{Date:} \hspace{1em} \rule{4cm}{0.4pt}
\vspace{2em}

\begin{tcolorbox}[colback=orange!30,colframe=black,title=Reminder]
This document is not used to reward AI use. It is used to judge your authorship, your decision-making, and the quality of your critical engagement. Aim for quality of reflection, not length. A specific account of one genuine interaction with AI is more useful to your marker than a vague account of many.
\end{tcolorbox}

\end{document}
